% Modernised reading — fonds Alexandre Grothendieck, shelfmark 13, pages 1-51.
%
% This is an INTERPRETATION, not a transcription. It re-reads the pages in
% today's notation and vocabulary, states the hypotheses the manuscript leaves
% implicit, and drops the critical apparatus so that the mathematics reads
% straight through. Everything the brackets and underlines carried moves into a
% footnote. It is held to being mathematically correct as it stands; where the
% manuscript is loose or mistaken, the true statement is what appears here and
% the footnote says what the page has. In any dispute of substance the
% transcription governs: cite batch-01/02/03.fr.tex.
%
% Pass: Opus 5 (claude-opus-5), 2026-09-19 — first pass, unchecked against the
% pages by a human.
%
% Scope: one document for the shelfmark, its three batches read as one
% argument. Forty-two of the fifty-one leaves carry mathematics; the rest are
% blank or are covers in his hand, whose words become section headings here.
%
% Spine. The folder asks ONE question, in five registers, and the question is:
% how much of a category of motives is recoverable from its fibre functors and
% the comparison isomorphisms between them? The recurring instrument is a
% group and a torsor — G = Aut(T_B) and P = Isom(T_B (x) k, T_DR) — with a
% transcendental point of P that is the period matrix.
%
%   pp. 3-11   THE HODGE GROUP. A functorial bigrading on Rep(G) is the same
%              thing as j : S -> G_R (S = Res_{C/R} G_m). Conditions (a), (b)
%              tie j to the central i : G_m -> G and to eps : G -> G_m. A
%              theorem characterises polarisability by the centraliser of
%              j(sqrt-1) being maximal compact. Then G-subordinations over an
%              analytic space, represented by the space J of such j, and the
%              moduli variety J/Gamma_0.
%   pp. 13-19  THE LETTER TO MUMFORD, Algiers 1 Nov. 1965, in English, and its
%              five-leaf postscript. He withdraws conjectures that contradict
%              Tate's; builds the same J from motivic Galois theory; observes
%              (margin, p. 16) that J is NOT in general riemannian symmetric;
%              and runs a Borel argument to an apparent contradiction with
%              Tate, closing « So I really do not know what to believe! ».
%   pp. 22-24  CONNECTIONS ON A FIBRE FUNCTOR. Connections form an affine space
%              over the canonical one, so one is an element omega(x) of g (x) K,
%              with a Maurer-Cartan integrability identity; the subspace n
%              generated by the omega(x) must be all of g for the criterion to
%              bite.
%   pp. 26-31  THE PERIOD TORSOR. P = Isom(T_B (x) k, T_DR) under G_k, the
%              point phi in P(C), the unipotent H' and the twisted G'', the
%              cocharacters i_1, i_2, the data a)-g), the 2 i pi as the identity
%              phi-dot = (1/2ipi) alpha_C, and the converse question answered by
%              « G'_C definable over k » — necessary, very restrictive, not
%              sufficient; to be completed by transcendence arguments.
%   pp. 33-34  DOUBLE LATTICES, the abstract engine: V (cap) T'(V) = V^H for a
%              group H built from a transporter, so full faithfulness is exactly
%              V^H = V^G.
%   pp. 36-43  THE ADELIC TABLE AND THE TRANSCENDENTAL DATUM: pi_1 -> G(A), the
%              cyclotomic character, Frobenius, the Tate modules; then
%              f : S-bar_C -> J holomorphic and equivariant, the monodromy group
%              R and the descent of (P/R, G/R) to k.
%   pp. 45-49  SIX CANDIDATE REALISATION FUNCTORS, each asked whether it is
%              fully faithful and under which conjecture; the « double niveau »
%              question reduced inside the Galois group to V^G = V (cap) g(V-bar).
%   pp. 50-51  A short closing note on characterising essentially algebraic de
%              Rham classes, ending on a question it does not answer.
%
% What only a folder-wide reading can say. Three links, none stated on any one
% leaf:
%   - the i_0 : U_R -> G_R of the letter (p. 15) IS the j' of p. 5: the letter's
%     data are the French run's data, re-derived four leaves later in English
%     and without the derivation. The letter is not an aside; it is the same
%     construction addressed to somebody else.
%   - the abstract « double lattice » criterion V^H = V^G of p. 34 IS the
%     answer-shape of the « double niveau » question of p. 47, where it
%     reappears as V^G = V (cap) g(V-bar) with H the group generated by the
%     transporter. Pages 33-34 prove in the abstract what page 47 then asks
%     about motives.
%   - the space he calls J on pp. 7-19 and Gamma_G on pp. 38-43 is one object
%     under two names, and the f : S-bar_C -> Gamma_G of p. 38 is the period map
%     of the G-subordinations of p. 7. The folder builds the same object twice,
%     eleven leaves apart, without remarking on it.
%
% NOTATION fixed for the folder. Collisions are his and are separated here:
%   G'   — pp. 3-11 he writes G' for ker(eps); pp. 26-43 for ad(P) =
%          Aut(T_DR). WRITTEN HERE: G^1 for ker(eps), G' for ad(P).
%   Q    — pp. 15-19 the flag variety G_R/C_R; pp. 27-43 the torsor
%          Isom_gr(T_DR, T_Hdg). WRITTEN HERE: the flag variety spelled out as
%          G_R/C_R, Q reserved for the torsor.
%   Gamma— p. 10 he sets Gamma = G(Q); pp. 38-43 Gamma = pi_1(S(C), xi_0), and
%          Gamma_G is the space of Hodge structures. WRITTEN HERE: Gamma is the
%          fundamental group, G(Q) is spelled out, and the space of Hodge
%          structures is J throughout.
%   P    — p. 10 the principal covering of S with group G(Q); pp. 26-43 the
%          period torsor. WRITTEN HERE: script P for the covering, P for the
%          torsor.
% S = Res_{C/R} G_m, U its norm-one subtorus, G_m the multiplicative group, A
% the finite adeles. T_B, T_DR, T_Hdg, T_{B-H} the fibre functors. g = Lie G.
%
% Substantive departures, each footnoted where it occurs:
%   - p. 5: the leaf says the conditions « sont équivalentes et impliquent
%     (iii) ». Under G reductive (iii) is EQUIVALENT to them, and is stated so
%     here; the implication is what the page writes.
%   - p. 5: reductivity of G is supplied only inside condition (i bis) on the
%     leaf. It is a hypothesis of the theorem and is stated as one here.
%   - pp. 26 ff: the Tannakian formalism the folder uses throughout presupposes
%     a neutral Tannakian category of motives — which is what the standard
%     conjectures would give and what the folder never establishes. Said once,
%     at the head of the period-torsor section, and not repeated.
%   - p. 27: the target of his (12) is written G' where G'' is meant; G'' is
%     what he has just identified with Aut(T_Hdg).
%   - p. 19: the question whether G(Z) is Zariski-dense when G_R is semisimple
%     without compact factors has an affirmative answer, by Borel's density
%     theorem. Footnoted; the page asks and does not answer.
%   - p. 34: he writes T'(W) where the argument carries T'(V). Restated in V.
%   - p. 38: two consecutive statements are both numbered 13). Renumbered here
%     13) and 14), with the discordance footnoted.
%   - p. 47: the containment is V^G (subset) V (cap) g(V-bar) and the question
%     is whether it is an equality. The direction is preserved; an equality
%     here would be the assertion, not the question.
%
% What the folder announces and never establishes: the standard conjectures
% (assumed wherever « moyennant Hodge » or « moyennant Tate » appears, which is
% most of it); the explicit form of conditions 2) and 3) of p. 28, of which he
% writes « On ne sait pas expliciter 2) et 3) directement en termes des données
% a) ... ! »; the analytic continuation formula for f and the relation between
% the transcendental datum D) and the adelic datum E) (p. 43); and the converse
% of condition b) of p. 50, which is the question the folder stops on.

\documentclass[11pt,a4paper]{article}
% Le préambule commun à toutes les transcriptions du fonds Grothendieck.
%
% Il tient en une idée : **l'incertitude se compose, elle ne se résout pas.**
% Une transcription qui lisse un mot douteux a détruit la seule chose pour
% laquelle on la faisait — un lecteur doit pouvoir savoir, à chaque endroit,
% ce qui était sur le feuillet et ce qui a été ajouté. Les six macros qui
% suivent sont donc l'appareil critique, et non de la décoration.
%
% Les mêmes commandes sont reconnues par `scripts/render.mjs`, qui produit la
% vue de lecture du site. Une macro ajoutée ici doit l'être aussi là-bas,
% faute de quoi le rendu échouera bruyamment — ce qui est voulu : un rendu
% silencieusement faux est pire qu'un rendu absent.

\ProvidesPackage{grothendieck}[2026/08/08 Transcriptions du fonds Grothendieck]

\usepackage{amsmath,amssymb,amsthm}
\usepackage{mathrsfs}
\usepackage{stmaryrd}
% Les diagrammes commutatifs sont partout dans ce fonds : c'est la langue
% même de la géométrie algébrique de Grothendieck, et une transcription qui
% les rendrait en prose ne transcrirait rien.
\usepackage{tikz-cd}
% Français par défaut — c'est la langue des feuillets — mais l'anglais est
% chargé aussi : la traduction et le résumé sont des documents anglais, et la
% typographie française (espaces avant les deux-points) y serait une faute.
% Ces documents basculent par \selectlanguage{english} en tête de corps.
\usepackage[english,french]{babel}
\usepackage{fontspec}
\usepackage{geometry}
\geometry{a4paper,margin=25mm,marginparwidth=28mm}
\usepackage{marginnote}
\usepackage{xcolor}
% ulem, not soul. `soul`'s \st and \ul are built on a character-by-character
% scan that XeLaTeX's font machinery breaks: the document fails with an
% undefined control sequence rather than degrading. ulem does the same two jobs
% and is Unicode-engine safe. `normalem` stops it hijacking \emph, which the
% transcriptions use for Grothendieck's own emphasis.
\usepackage[normalem]{ulem}
\usepackage{hyperref}
\hypersetup{colorlinks=true,linkcolor=black,urlcolor=black,pdfborder={0 0 0}}

% --- Métadonnées du lot -----------------------------------------------------
% Reprises telles quelles par le script de rendu, qui les affiche en tête de
% la vue de lecture. Elles fixent ce que le fichier prétend couvrir : sans
% elles, une transcription flotte sans point d'attache dans un fonds de seize
% mille feuillets.

\newcommand{\folder}[1]{\gdef\grothFolder{#1}}
\newcommand{\batch}[1]{\gdef\grothBatch{#1}}
\newcommand{\leaves}[2]{\gdef\grothFirst{#1}\gdef\grothLast{#2}}
\newcommand{\foldertitle}[1]{\gdef\grothFoldertitle{#1}}
\gdef\grothFolder{?}\gdef\grothBatch{?}\gdef\grothFirst{?}\gdef\grothLast{?}\gdef\grothFoldertitle{}

% --- Le résumé --------------------------------------------------------------
%
% Il ouvre la lecture modernisée, sous le titre « Résumé », et il est composé
% en petit. Ce n'est pas un ornement : c'est le seul passage du document qui
% ne soit pas des mathématiques, et un lecteur doit pouvoir le reconnaître
% avant de l'avoir lu — puis le sauter s'il connaît déjà le sujet, ou s'y
% arrêter s'il ne le connaît pas. Le filet qui le suit marque la reprise du
% corps.

\newenvironment{resume}
  {\small\setlength{\parskip}{0.45em}}
  {\par\medskip\hrule\medskip}

% --- Mots-clés ---------------------------------------------------------------
%
% En anglais, en clôture du résumé de la lecture modernisée : le vocabulaire
% moderne sous lequel chercher ce que les feuillets construisent. Le manifeste
% du site (`npm run manifest`) les extrait pour étiqueter la cote entière —
% cette ligne est la source unique des tags, il n'y a pas de fichier à côté.
\newcommand{\keywords}[1]{\par\smallskip\noindent
  {\footnotesize\textsc{Keywords} --- \textit{#1}}\par}

% --- L'appareil critique ----------------------------------------------------

% Le numéro de feuillet, en marge. C'est l'ancre qui permet de régler un
% désaccord sur une formule en regardant *un* feuillet plutôt qu'une chemise
% de six cents. Le site s'en sert aussi pour tourner les pages du fac-similé
% quand on fait défiler la transcription.
\newcommand{\leaf}[1]{%
  \marginnote{\footnotesize\color{gray!70}\textsf{#1}}%
  \label{leaf:#1}\ignorespaces}

% Illisible : on ne devine pas. Un mot inventé qui se lit comme les autres est
% le pire résultat possible.
\newcommand{\ill}{\textcolor{red!60!black}{[\,\ldots\,]}}

% Lecture douteuse : le mot est donné, et signalé comme incertain.
\newcommand{\uncertain}[1]{\uline{#1}}

% Ajout de l'éditeur — une lettre manquante, un mot sous-entendu.
\newcommand{\add}[1]{\textcolor{blue!50!black}{[#1]}}

% Biffé par Grothendieck lui-même. Ce qu'il a rayé fait partie du document :
% une rature dit souvent où la pensée a changé de direction.
\newcommand{\struck}[1]{\sout{#1}}

% Note du transcripteur, sur l'état matériel du feuillet.
\newcommand{\note}[1]{\par\smallskip{\small\color{gray!60!black}\textit{[#1]}}\par\smallskip}

% Note marginale de Grothendieck, distincte du corps du texte.
\newcommand{\marginal}[1]{\par\smallskip\noindent\hspace{1em}%
  {\small\color{gray!60!black}#1}\par\smallskip}

% --- Titre ------------------------------------------------------------------

\AtBeginDocument{%
  \begin{center}
    {\large\bfseries Fonds Alexandre Grothendieck --- cote n\textsuperscript{o}~\grothFolder}\\[2pt]
    {\itshape\grothFoldertitle}\\[4pt]
    {\small feuillets \grothFirst--\grothLast\ (lot \grothBatch)}\\[2pt]
    {\footnotesize\color{gray!60!black}Transcription établie d'après le fac-similé numérisé
     par l'Université de Montpellier.\\
     Les lectures incertaines sont soulignées, les illisibles notées [\,\ldots\,],
     les ajouts entre crochets.}
  \end{center}
  \vspace{1em}}

\endinput


\folder{13}
\pages{1}{51}
\dating{1965}
\watermark{Édition de démonstration\\interprétation personnelle de l'œuvre}
\foldertitle{Groupe de Galois motivique : notes manuscrites (s.d.), lettre (1965). — lecture modernisée du dossier entier}

\begin{document}

\section*{Résumé}

\begin{resume}
Ces feuillets demandent ce qu'il reste d'une variété algébrique quand on ne
garde d'elle que ses cohomologies et les dictionnaires qui les comparent.

Voici de quoi il s'agit. À une variété algébrique --- disons, pour fixer les
idées, une courbe ou une surface définie par des équations à coefficients
rationnels --- on sait associer plusieurs espaces vectoriels qui en mesurent la
forme. Le premier est topologique : on oublie les équations, on ne regarde que
l'objet géométrique, et l'on compte ses trous. C'est la cohomologie
\emph{de Betti}, et elle vit sur les nombres rationnels. Le second est
différentiel : on considère les formes différentielles que l'on peut écrire sur
la variété, modulo celles qui sont des différentielles de quelque chose. C'est
la cohomologie \emph{de De Rham}, et elle vit sur le corps des coefficients des
équations. Le troisième est arithmétique : on regarde comment le groupe de
Galois du corps de base agit sur les points de torsion. Ces trois espaces ont
la même dimension, et ils ne vivent pas sur le même corps.

Le problème naît du dictionnaire entre les deux premiers. Il existe, et c'est
un théorème ; mais pour l'écrire il faut sortir des nombres rationnels. Le cas
le plus simple est celui du cercle privé d'un point : la forme différentielle
$dz/z$ engendre la cohomologie de De Rham, le tour complet engendre la
cohomologie de Betti, et le dictionnaire entre les deux est l'intégrale de
l'une le long de l'autre, qui vaut $2i\pi$. Ce nombre n'est pas rationnel, et
n'est même pas algébrique. Dans le cas général le dictionnaire est une matrice,
dont les coefficients s'appellent les \emph{périodes} de la variété, et dont
$2i\pi$ est l'entrée la plus modeste.

L'idée qui organise tout le dossier est celle-ci. On ne sait pas grand-chose
d'une matrice de périodes prise isolément ; on sait en revanche parler de la
\emph{position} d'une structure rationnelle par rapport à une autre à
l'intérieur d'un même espace complexe. Prenons un espace vectoriel complexe et
deux réseaux rationnels dedans. Le passage de l'un à l'autre est une matrice de
changement de base, mais cette matrice n'est définie qu'à la symétrie de chacun
des deux réseaux près : ce qui est intrinsèque, c'est une classe, un point de ce
que l'on appelle un \emph{torseur}. Un torseur, c'est un espace sur lequel un
groupe agit simplement transitivement --- comme le plan sur lequel agissent les
translations : aucun point n'y est privilégié, mais deux points quelconques
diffèrent d'un unique élément du groupe. Il n'a pas d'origine, et c'est
précisément pour cela qu'il est le bon objet : une matrice de périodes n'a pas
d'origine non plus.

Le groupe qui agit est le \emph{groupe de Galois motivique}. On l'obtient en
prenant toutes les symétries de la cohomologie de Betti qui respectent tout ce
qui est algébrique --- les produits, les classes de cycles, bref tout ce que la
géométrie fournit. Il joue, pour les cohomologies d'une variété, le rôle que le
groupe de Galois ordinaire joue pour les racines d'un polynôme : plus il est
gros, moins il y a de relations, et réciproquement. Le torseur des périodes est
alors l'espace des dictionnaires possibles entre Betti et De Rham, et la matrice
des périodes en est un point complexe. Tout le dossier consiste à faire
travailler ce couple --- un groupe, et un torseur sous ce groupe avec un point
transcendant.

Une seconde idée se superpose à la première. Sur la cohomologie d'une variété
complexe, les nombres complexes agissent, d'une manière qui enregistre le
découpage des formes différentielles selon le nombre de $dz$ et de
$d\bar{z}$ qu'elles portent. Se donner ce découpage, c'est se donner une action
du groupe des nombres complexes non nuls --- vu comme groupe algébrique réel,
il porte ici le nom de $\mathbb{S}$. Une \emph{structure de Hodge} n'est rien
d'autre qu'une telle action ; et si l'on a fixé le groupe $G$, l'ensemble des
structures de Hodge possibles est l'ensemble des morphismes de $\mathbb{S}$
dans $G$, qui forme un espace noté ici $\mathcal{J}$. Lorsque la variété varie
dans une famille, sa structure de Hodge varie avec elle, et l'on obtient une
application de l'espace des paramètres vers $\mathcal{J}$ : c'est l'\emph{application
des périodes}. Les quotients $\mathcal{J}/\Gamma$ par un groupe arithmétique
sont les « variétés modulaires » que le dossier cherche à reconnaître.

C'est ici que le dossier devient inquiet, et c'est ce qui en fait le prix. Une
lettre à Mumford, écrite d'Alger le 1er novembre 1965, ouvre en retirant des
conjectures qu'il avait avancées : sous la forme où il les avait énoncées, elles
contredisent celles de Tate. Suit un post-scriptum de cinq feuillets où il
construit l'espace $\mathcal{J}$, observe en marge qu'il n'est pas en général de
la forme la plus favorable, et mène un argument qui aboutit à une contradiction
apparente avec Tate. La lettre se termine sur six mots : \emph{So I really do not
know what to believe!} Le reste du dossier travaille dans cette incertitude, et
les trois dernières lignes du dernier feuillet posent une question de plus, à
laquelle rien ne répond : le dossier ne se termine pas, il s'arrête.

Les noms sous lesquels on cherchera la suite : groupe de Mumford--Tate et
domaine de Mumford--Tate, variétés de Shimura, torseur des périodes et
conjecture des périodes, variations de structures de Hodge et application des
périodes, conjectures de Hodge et de Tate, et pleine fidélité des foncteurs de
réalisation.

\keywords{motivic Galois group, period torsor, Mumford-Tate domain, Shimura datum, Deligne torus, Cartan involution, polarisable Hodge structure, period map, variation of Hodge structure, monodromy group, theorem of the fixed part, Tannakian formalism, fibre functor, integrable connection on a fibre functor, Gauss-Manin stratification, full faithfulness of realisation functors, Hodge conjecture, Tate conjecture, Griffiths transversality, cyclotomic character, adelic representation, Frobenius element, l-adic monodromy, Borel density theorem, Siegel upper half-space, arithmetic group, canonical model, double lattice, transcendence of periods}
\end{resume}

\pagerange{3}{3}
\section*{Le fil du dossier, et les conventions}

\emph{Les stations.}

\begin{enumerate}
  \item \emph{Pages 3 à 11}, « Groupe de Hodge » : se donner, de façon
  fonctorielle, une structure de Hodge sur toutes les représentations d'un
  groupe $G$, c'est se donner un morphisme $j : \mathbb{S} \to G_{\mathbb{R}}$.
  Un théorème caractérise ceux pour lesquels ces structures sont polarisables.
  L'espace $\mathcal{J}$ de ces $j$ représente un foncteur, et son quotient par
  un groupe arithmétique est une variété de modules.
  \item \emph{Pages 13 à 19}, « Lettre à Mumford » : la même construction,
  refaite en anglais et sans les démonstrations, suivie d'un argument qui bute
  sur les conjectures de Tate.
  \item \emph{Pages 22 à 24} : la situation abstraite d'une connexion
  intégrable sur un foncteur fibre. C'est l'outil technique dont les pages 38 à
  43 se serviront.
  \item \emph{Pages 26 à 31} : le torseur des périodes, la donnée $a)$ à $g)$
  qui code une catégorie de motifs, le facteur $2i\pi$, et la question inverse.
  \item \emph{Pages 33 et 34} : les « doubles réseaux », c'est-à-dire le
  critère abstrait de pleine fidélité $V^{H} = V^{G}$.
  \item \emph{Pages 36 et 37} : la récapitulation adélique, sept encadrés et
  neuf propriétés.
  \item \emph{Pages 38 à 43} : la donnée transcendante, le revêtement
  équivariant, le groupe de monodromie et la descente à $k$.
  \item \emph{Pages 45 à 49} : six foncteurs de réalisation candidats, chacun
  interrogé sur sa pleine fidélité et sur la conjecture qu'elle coûte.
  \item \emph{Pages 50 et 51} : les classes de De Rham essentiellement
  algébriques, et la question sur laquelle le dossier s'arrête.
\end{enumerate}

\emph{Ce que seule une lecture d'ensemble peut dire.} Trois liens, qu'aucun
feuillet n'énonce.

Le premier : le morphisme $i_{0} : \mathbb{U}_{\mathbb{R}} \to G_{\mathbb{R}}$
par lequel s'ouvre le post-scriptum de la lettre (page 15) est exactement le
$j'$ de la page 5. La lettre ne fait pas une digression ; elle refait la même
construction, quatre feuillets plus loin, en anglais et sans la démonstration,
pour quelqu'un d'autre.

Le deuxième : le critère abstrait $V \cap T'(V) = V^{H}$ établi aux pages 33 et
34 est la forme même de la réponse que la page 47 cherche, où il reparaît sous
les espèces de $V^{G} = V \cap g(\overline{V})$. Les pages 33 et 34 démontrent
dans l'abstrait ce que la page 47 demandera des motifs.

Le troisième : l'espace qu'il note $\mathcal{J}$ aux pages 7 à 19 et
$\Gamma_{G}$ aux pages 38 à 43 est un seul objet sous deux noms, et
l'application $f$ de la page 38 est l'application des périodes des
$G$-subordinations de la page 7. Le dossier construit deux fois le même objet, à
onze feuillets de distance, sans le remarquer.

\emph{Les conventions.} Quatre lettres servent chacune à deux choses dans le
dossier, et sont séparées ici.

\begin{itemize}
  \item $G'$ désigne, aux pages 3 à 11, le noyau de $\varepsilon$, et à partir
  de la page 26 le groupe tordu $\operatorname{ad}(P) \simeq
  \operatorname{Aut}^{\otimes}(T_{DR})$. On écrit ici $G^{1}$ pour le noyau de
  $\varepsilon$, et l'on réserve $G'$ au second.
  \item $Q$ désigne, aux pages 15 à 19, la variété de drapeaux
  $G_{\mathbb{R}}/C_{\mathbb{R}}$, et à partir de la page 27 le torseur
  $\operatorname{Isom}^{\otimes}(T_{DR}, T_{\mathrm{Hdg}})$. On écrit ici la
  variété de drapeaux en toutes lettres, et l'on réserve $Q$ au torseur.
  \item $\Gamma$ vaut $G(\mathbb{Q})$ à la page 10 et
  $\pi_{1}(S(\mathbb{C}), \xi_{0})$ aux pages 38 à 43, cependant que
  $\Gamma_{G}$ est un espace de structures de Hodge. On réserve ici $\Gamma$ au
  groupe fondamental, on écrit $G(\mathbb{Q})$ en toutes lettres, et l'espace
  des structures de Hodge s'appelle $\mathcal{J}$ d'un bout à l'autre.
  \item $P$ désigne à la page 10 le revêtement principal de $S$ de groupe
  $G(\mathbb{Q})$, et à partir de la page 26 le torseur des périodes. On écrit
  ici $\mathcal{P}$ le revêtement et $P$ le torseur.
\end{itemize}

Pour le reste, $\mathbb{S} = \prod_{\mathbb{C}/\mathbb{R}} \mathbb{G}_{m}$ est
le tore de Weil des restrictions des scalaires\footnote{Ce tore, dont le groupe
des points réels est $\mathbb{C}^{*}$, est aujourd'hui couramment appelé
\emph{tore de Deligne} ; le nom est postérieur aux feuillets, qui l'emploient
sans le nommer.}, $\mathbb{U} \subset \mathbb{S}$ son sous-tore des éléments de
norme $1$, $\mathbb{G}_{m}$ le groupe multiplicatif, $\mathbb{A}$ l'anneau des
adèles finis de $\mathbb{Q}$, et $\mathfrak{g} = \operatorname{Lie} G$. Les
foncteurs fibres sont $T_{B}$ (Betti), $T_{DR}$ (De Rham),
$T_{\mathrm{Hdg}} = \operatorname{gr}^{*} T_{DR}$ et $T_{B-H}$ (Betti muni de sa
structure de Hodge).

\emph{Ce que le dossier suppose sans l'établir.} Les conjectures standard, dont
il a besoin partout où il écrit « moyennant Hodge » ou « moyennant Tate »,
c'est-à-dire presque partout ; la forme explicite des conditions $2)$ et $3)$ de
la page 28, dont il note lui-même qu'on ne sait pas l'écrire ; la formule de
prolongement analytique de $f$ et la relation entre les données $D)$ et $E)$ de
la page 43 ; et la réciproque de la condition $b)$ de la page 50, qui est la
question sur laquelle il s'arrête.

\pagerange{3}{11}
\section*{Le groupe de Hodge (pages 3 à 11)}

\subsection*{Une bigraduation fonctorielle est un morphisme de tores}

Soit $\mathbb{S} = \prod_{\mathbb{C}/\mathbb{R}} \mathbb{G}_{m,\mathbb{C}}$,
de sorte que $\mathbb{S}(\mathbb{R}) \simeq \mathbb{C}^{*}$, et soit
$\mathbb{U} = \operatorname{Ker}(\text{Norme} : \mathbb{S} \to
\mathbb{G}_{m,\mathbb{R}})$. La catégorie $\operatorname{Rep}(\mathbb{S})$
s'identifie à celle des espaces vectoriels réels $V_{\mathbb{R}}$ munis d'une
bigraduation
\[
V_{\mathbb{C}} \;=\; \bigoplus_{p,q} V^{p,q} , \qquad V^{p,q} = \overline{V^{q,p}} ,
\]
l'élément $\lambda \in \mathbb{C}^{*} \simeq \mathbb{S}(\mathbb{R})$ opérant sur
$V^{p,q}$ par $\lambda^{p} \bar{\lambda}^{q}$. Le morphisme canonique
$i : \mathbb{G}_{m,\mathbb{R}} \to \mathbb{S}$ correspond à la graduation
totale.\footnote{Le feuillet écrit ces sommes directes avec le signe $\coprod$,
et s'y tient ; on a rétabli $\bigoplus$, qui est ce que l'énoncé veut dire.}

Soit maintenant $G$ un groupe algébrique sur $\mathbb{Q}$. Se donner, sur le
foncteur composé
\[
\operatorname{Rep}(G) \longrightarrow \operatorname{Mod}_{f}(\mathbb{Q})
\longrightarrow \operatorname{Mod}_{f}(\mathbb{R})
\longrightarrow \operatorname{Mod}_{f}(\mathbb{C}) ,
\]
une bigraduation fonctorielle vérifiant $V^{p,q} = \overline{V^{q,p}}$ ---
c'est-à-dire une préstructure de Hodge compatible avec la structure rationnelle
et fonctorielle en l'objet --- revient à se donner un morphisme
\[
j : \mathbb{S} \longrightarrow G_{\mathbb{R}} .
\]

On se donne de plus un morphisme central $i : \mathbb{G}_{m} \to G$, c'est-à-dire
une graduation $V = \bigoplus_{n} V^{n}$ du foncteur oubli, et un objet
inversible $\mathbb{Q}(-1)$ de $\operatorname{Rep}(G)$, c'est-à-dire un morphisme
$\varepsilon : G \to \mathbb{G}_{m}$ soumis à $\varepsilon\, i(\lambda) =
\lambda^{2}$. Deux conditions relient $j$ à ces données :
\[
(\mathrm{a}) \qquad j\big|\,\mathbb{G}_{m,\mathbb{R}} = i_{\mathbb{R}} ,
\]
qui exprime que la graduation totale est bien celle donnée par $i$, autrement dit
que $V^{n}_{\mathbb{C}} = \bigoplus_{p+q=n} V^{p,q}$ ; et
\[
(\mathrm{b}) \qquad \varepsilon\, j = \text{Norme} : \mathbb{S} \longrightarrow \mathbb{G}_{m,\mathbb{R}} ,
\]
qui exprime que $\mathbb{Q}(1)$ est de bidegré $(-1,-1)$.\footnote{Le feuillet 4
écrit que $\mathbb{Q}(1)$ est « de bidegré $(1,1)$ » et développe la condition
$(\mathrm{b})$ sous la forme $\varepsilon j(\lambda) = \lambda^{2}$ pour
$\lambda$ réel et $\varepsilon(j(u)) = 1$ pour $u \in \mathbb{U}$, qui est la
bonne. Le signe du bidegré dépend de la convention choisie pour $\mathbb{Q}(1)$
et de celle prise pour $\varepsilon$ ; la condition $(\mathrm{b})$, elle, ne
dépend d'aucune des deux et c'est elle qui porte.}

\subsection*{Quand les structures obtenues sont-elles polarisables}

Fixons $i$ et $\varepsilon$, et considérons les $j$ vérifiant $(\mathrm{a})$ et
$(\mathrm{b})$. Pour $V$ homogène de degré $n$, une forme
$\varphi : V \otimes V \to \mathbb{Q}(-n)$ de $\operatorname{Rep}(G)$ polarise la
structure de Hodge associée à $j$ lorsque
\[
(x,y) \longmapsto \varphi_{\mathbb{R}}(x, Cy)
\]
est définie positive sur $V_{\mathbb{R}}$, où $C = j(\sqrt{-1})$. Une telle forme
est nécessairement symétrique si $n$ est pair, alternée si $n$ est impair.

Posons $G^{1} = \operatorname{Ker} \varepsilon$. On a alors
$G \simeq G^{1} \times_{\mathbb{Z}/2\mathbb{Z}} \mathbb{G}_{m,\mathbb{Q}}$, où
$\mathbb{Z}/2\mathbb{Z} \to G^{1}$ est induit par $i$ : le triplet
$(G, i, \varepsilon)$ se reconstitue donc à partir de $G^{1}$ et de son point
d'ordre $2$, $\eta = i(-1)$. De même, se donner $j : \mathbb{S} \to
G_{\mathbb{R}}$ vérifiant $(\mathrm{a})$ et $(\mathrm{b})$ équivaut à se donner
$j' : \mathbb{U} \to G^{1}_{\mathbb{R}}$ vérifiant $j'(-1) = \eta$.

\textbf{Théorème.} \emph{Soient $G$ réductif sur $\mathbb{Q}$, $i$ central et
$\varepsilon$ comme ci-dessus, et $j : \mathbb{S} \to G_{\mathbb{R}}$ vérifiant
$(\mathrm{a})$ et $(\mathrm{b})$. Les conditions suivantes sont
équivalentes.}\footnote{Le feuillet 5 écrit que les conditions « sont
équivalentes et impliquent (iii) ». Sous l'hypothèse de réductivité, (iii) leur
est en fait équivalente, et c'est ainsi qu'on l'énonce ici ; l'implication seule
est ce que porte la page. La réductivité, de son côté, n'est supposée sur le
feuillet qu'à l'intérieur de la condition (i bis) ; c'est une hypothèse du
théorème et elle est placée ici comme telle.}

\begin{enumerate}
  \item[(i)] Pour tout $V \in \operatorname{Rep}(G)$ homogène de degré $n$, il
  existe dans $\operatorname{Rep}(G)$ une forme
  $\varphi : V \otimes V \to \mathbb{Q}(-n)$ qui polarise la structure de Hodge
  définie par $j$.
  \item[(i bis)] La même condition, restreinte à une famille $(V_{\alpha})$
  telle que $G \to \prod_{\alpha} \operatorname{Aut}_{\mathbb{Q}}(V_{\alpha})$
  soit injectif.
  \item[(ii)] Les mêmes conditions, avec $G$ remplacé par $G_{\mathbb{R}}$ et
  $\operatorname{Rep}(G)$ par $\operatorname{Rep}(G_{\mathbb{R}})$.
  \item[(iii)] Le centralisateur de $C = j(\sqrt{-1})$ dans
  $(G^{1})^{0}_{\mathbb{R}}$ est un sous-groupe compact maximal.
\end{enumerate}

La condition $(\mathrm{iii})$ est la formulation, dans le langage des groupes,
de ce que $\operatorname{ad}(C)$ est une involution de Cartan de
$G^{\mathrm{ad}}_{\mathbb{R}}$ : c'est bien une involution, car $C^{2} = j(-1) =
i(-1)$ est central par $(\mathrm{a})$, et son groupe des points fixes est
exactement le centralisateur de $C$.\footnote{On reconnaît le critère de
polarisabilité tel qu'on l'énonce aujourd'hui pour les données de Shimura ;
l'axiomatisation sous ce nom est postérieure aux feuillets, qui n'en disposent
pas.}

\subsection*{$G$-subordinations, et l'espace qui les classe}

Soit $S$ un espace analytique. Une \emph{$G$-subordination} sur $S$ est la
donnée d'un $\otimes$-foncteur qui, à tout $V \in \operatorname{Rep}(G)$, associe
une structure de Hodge sur le fibré $V_{S}$ --- c'est-à-dire une bigraduation
analytique réelle $V_{S}^{p,q} = \overline{V_{S}^{q,p}}$ dont la filtration
associée est analytique complexe --- de manière que le degré total soit celui
défini par $i$, que $\mathbb{Q}(-1)$ soit de bidegré $(1,1)$, et que la fibre en
chaque point soit polarisable par une forme $G$-invariante. Lorsque $S$ est
connexe, il suffit d'exiger la polarisabilité en un seul point.

Traduites, ces conditions disent exactement ceci : pour chaque $s \in S$ on veut
un $j = J(s) : \mathbb{S} \to G_{\mathbb{R}}$ satisfaisant les conditions du
théorème, c'est-à-dire un $j' = J'(s) : \mathbb{U} \to G^{1}_{\mathbb{R}}$ avec
$j'(-1) = \eta$ et centralisateur de $j'(\sqrt{-1})$ compact maximal dans
$(G^{1})^{0}_{\mathbb{R}}$ ; et l'on veut que $J'(s)$ dépende de $s$ de façon
analytique réelle.

Soit alors $\mathcal{J}$ l'ensemble des $j$ (ou, ce qui revient au même, des
$j'$) possibles. Pour exprimer la condition d'holomorphie il convient de munir
$\mathcal{J}$ d'une structure de variété analytique complexe naturelle, et la
condition devient l'holomorphie de $s \mapsto j(s)$. On obtient ainsi :

\begin{quote}
L'ensemble des $G$-subordinations rigidifiées sur $S$ s'identifie à
$\operatorname{Mor}(S, \mathcal{J})$. Le foncteur $S \mapsto \{\,
G\text{-subordinations rigidifiées sur } S \,\}$ est donc représentable, et
$\mathcal{J}$ le représente.
\end{quote}

Le groupe $G(\mathbb{R})$ opère sur $\mathcal{J}$, avec des orbites ouvertes, et
les composantes connexes de $\mathcal{J}$ sont de la forme
\[
G(\mathbb{R})^{\circ} \big/ H(\mathbb{R}) \cap G(\mathbb{R})^{\circ}
\qquad\text{ou encore}\qquad
G^{1}(\mathbb{R})^{\circ} \big/ H'(\mathbb{R}) \cap G^{1}(\mathbb{R})^{\circ} ,
\]
où $H$ est le centralisateur d'un $j$ dans $G$, et $H'$ celui d'un $j'$ dans
$G^{1}$.\footnote{C'est ce qu'on appelle aujourd'hui un \emph{domaine de
Mumford--Tate} : l'orbite, sous $G(\mathbb{R})^{\circ}$, d'un morphisme
$j$ satisfaisant les conditions du théorème. La page 16 observera qu'il n'est pas
en général hermitien symétrique, ce qui est précisément la raison pour laquelle
ces domaines ne donnent pas toujours des variétés de Shimura.}

\subsection*{Sans point-base : le torseur et le revêtement}

La présentation précédente choisit un point-base. En voici une qui n'en choisit
pas. Soit $S$ connexe, $\mathcal{C} = \operatorname{Rep}(G)$, et
$F : \mathcal{C} \to \operatorname{Hodge}(S)$ un $\otimes$-foncteur muni d'un
isomorphisme $\varphi$ entre le foncteur sous-jacent et le foncteur oubli
$\omega$. Le groupe $\pi = \pi_{1}(S, \xi)$ opère sur le foncteur « fibrés
rationnels sur $S$ », d'où sur le composé avec $F$, d'où un morphisme
\[
\alpha : \pi_{1}(S, \xi) \longrightarrow G(\mathbb{Q}) .
\]
D'autre part, si $\widetilde{S}$ désigne le revêtement universel de $S$ pointé
au-dessus de $s$, la donnée de $\varphi$ permet de voir $V \mapsto
F(V)|_{\widetilde{S}}$ comme une $G$-subordination sur $\widetilde{S}$, d'où un
second élément de structure
\[
f : \widetilde{S} \longrightarrow \mathcal{J} ,
\]
les deux étant reliés par la compatibilité de $f$ aux actions de $\pi$, via
$\alpha$ et l'action de $G(\mathbb{R})$ sur $\mathcal{J}$. Se donner un couple
$(F, \varphi)$ équivaut à se donner un couple $(f, \alpha)$ soumis à ces
conditions.

Voici enfin la forme intrinsèque. Pour $s \in S$, soit $\Phi_{s}$ le foncteur
fibre en $s$ sur $\operatorname{Hodge}(S)$, et
\[
P_{s} \;=\; \operatorname{Isom}^{\otimes}(\omega,\ \Phi_{s} F) ,
\]
torseur sous $G(\mathbb{Q})$. Quand $s$ varie, les $P_{s}$ forment un revêtement
principal $\mathcal{P}$ de $S$ de groupe $G(\mathbb{Q})$, et le composé
\[
\operatorname{Rep}(G) \xrightarrow{\ F\ } \operatorname{Hodge}(S)
\xrightarrow{\ V\, \mapsto\, V \times_{S} \mathcal{P}\ } \operatorname{Hodge}(\mathcal{P})
\]
définit sur $\mathcal{P}$ une $G$-subordination, donc un morphisme
$g : \mathcal{P} \to \mathcal{J}$ compatible à l'action de $G(\mathbb{Q})$.
La catégorie des $G$-subordinations sur $S$ est alors équivalente à celle des
couples $(\mathcal{P}, g)$ formés d'un revêtement principal de groupe
$G(\mathbb{Q})$ et d'un morphisme équivariant vers $\mathcal{J}$.

Soit enfin $\Gamma_{0} \subset G(\mathbb{Q})$ un sous-groupe commensurable à un
$G(\mathbb{Z})$. Une \emph{$G$-subordination $\Gamma_{0}$-rigidifiée} sur $S$
est la donnée d'un revêtement principal $\mathcal{P}_{0}$ de groupe
$\Gamma_{0}$ et d'un $\Gamma_{0}$-morphisme $\mathcal{P}_{0} \to \mathcal{J}$ ;
pour $\Gamma_{0} = \{e\}$ on retrouve les subordinations rigidifiées, pour
$\Gamma_{0} = G(\mathbb{Q})$ les subordinations tout court. Cette structure
admet un espace de modules, à savoir
\[
\mathcal{J} / \Gamma_{0} ,
\]
qui porte une structure d'espace analytique dès que $\Gamma_{0}$ opère
proprement et discontinûment : c'est la \emph{variété de modules de niveau
$\Gamma_{0}$}.\footnote{Le feuillet 11 s'arrête là, sa moitié basse blanche.
Rien de cette construction ne passe dans la liasse suivante ; elle reparaît en
revanche, sous d'autres noms, aux pages 15 à 19 puis 38 à 43.}

\pagerange{13}{19}
\section*{Lettre à Mumford (pages 13 à 19)}

\subsection*{La lettre}

Elle est datée d'Alger, le 1er novembre 1965, et elle est en anglais ; elle
le reste dans la transcription et on n'en traduit ici que la teneur.

Il commence par retirer ce qu'il avait avancé. Les conjectures sur les
structures de Hodge dont il avait parlé à Mumford contredisent, sous la forme
où il les avait énoncées, celles de Tate. Celles-ci impliquent en effet ceci :
si $X$ est un schéma lisse, connexe et simplement connexe sur $\mathbb{C}$, et
$V$ une famille analytique complexe polarisée de structures de Hodge
paramétrée par $X^{\mathrm{hol}}$, alors $V$ n'est « algébrique » que si elle
est constante. Or il est facile de produire des familles non constantes, en
prenant pour $X$ un espace homogène sous un groupe affine sur $\mathbb{C}$ ;
elles ne devraient donc pas être algébriques, si l'on veut que les conjectures
de Tate tiennent.

Il propose alors un test : peut-on vraiment ne pas obtenir de $V$ algébrique de
cette manière ? Et il pose une question concrète, qui est la seule chose de
toute la lettre qu'on puisse chercher à vérifier sur un exemple :

\begin{quote}
Connaît-on une surface projective lisse $Y$ sur $\mathbb{C}$ telle que
$H^{1,1}(Y, \mathbb{C})$ soit engendré par des cycles algébriques et que
$\dim H^{0,2}(Y, \mathbb{C}) \geqslant 2$ ?
\end{quote}

\subsection*{Le post-scriptum : les variétés modulaires de la théorie de Galois motivique}

Suivent cinq feuillets qui refont, en anglais et sans les démonstrations, la
construction des pages 3 à 11.

Soit $G$ un groupe algébrique réductif connexe sur $\mathbb{Q}$, et
$\mathbb{U}_{\mathbb{R}}$ le tore réel tordu de dimension $1$, dont le groupe
des points réels est celui des nombres complexes de module $1$. On se donne
\[
i_{0} : \mathbb{U}_{\mathbb{R}} \longrightarrow G_{\mathbb{R}}
\]
tel que $i_{0}(-1)$ soit central, que le centralisateur $K_{\mathbb{R}}$ de
$i_{0}(\sqrt{-1})$ vérifie que $K_{\mathbb{R}}(\mathbb{R})^{\circ}$ est compact
maximal dans $G_{\mathbb{R}}(\mathbb{R})^{\circ}$, et que la composante de
$i_{0}$ soit non triviale pour chaque facteur $\mathbb{R}$-simple de
$G_{\mathbb{R}}$. Ce $i_{0}$ est exactement le $j'$ de la page 5, et les deux
premières conditions sont mot pour mot celles du théorème ; la troisième est
nouvelle et écarte les facteurs sur lesquels rien ne varierait.

Posant $C_{\mathbb{R}} = \operatorname{Centr}_{G_{\mathbb{R}}}(i_{0})$, et
notant $G_{\mathbb{R}}/C_{\mathbb{R}}$ la variété de drapeaux
correspondante,\footnote{Le feuillet note cette variété $Q$ ; c'est une autre
lettre que le corps des rationnels, et que le torseur $Q$ des pages 27 et
suivantes. On l'écrit ici en toutes lettres.} il appelle « Siegel-space »
\[
\mathcal{J} \;=\; G(\mathbb{R})^{\circ} \big/ G(\mathbb{R})^{\circ} \cap C_{\mathbb{R}}(\mathbb{R})
\;\subset\; (G_{\mathbb{R}}/C_{\mathbb{R}})(\mathbb{R}) .
\]
C'est le même $\mathcal{J}$ qu'à la page 7, et il hérite de la variété de
drapeaux une structure complexe naturelle, invariante sous $G(\mathbb{R})$.
Toute représentation rationnelle de $G$ dans un $V$ sur $\mathbb{Q}$, munie d'un
réseau $V_{0}$, définit alors sur $\mathcal{J}$ une famille analytique complexe
de structures de Hodge de degré total $n$ fixé --- nécessairement globalement
polarisable, une polarisation s'obtenant en choisissant sur $V$ une forme
invariante sous $G$ --- pourvu que $i_{0}(-1)$ opère sur $V$ par l'homothétie
$(-1)^{n}$.

Vient alors, en marge, l'observation qui décide de tout ce qui suit :

\begin{quote}
En général $\mathcal{J}$ n'est \emph{pas} riemannien symétrique, mais est un
espace fibré au-dessus d'un espace riemannien symétrique, de fibres des espaces
homogènes projectifs complexes sous $K_{\mathbb{C}}$. Cela implique au moins que
$\mathcal{J}$ est \emph{simplement connexe}, donc est le revêtement universel de
$M = \mathcal{J}/\Gamma$ lorsque $\Gamma$ opère librement.
\end{quote}

C'est là le point où le domaine cesse d'être un domaine hermitien symétrique et
où les quotients cessent d'être automatiquement des variétés algébriques ---
c'est-à-dire, en termes d'aujourd'hui, le point où un domaine de Mumford--Tate
cesse d'être une donnée de Shimura. Que $\mathcal{J}$ reste simplement connexe
est ce qui sauve le quotient comme espace analytique et ce qui fera marcher, à
la page 18, l'argument de Borel.

Si $\Gamma \subset G(\mathbb{Q})$ est le groupe laissant $V_{0}$ fixe, il opère
sur $(V, \mathcal{J})$, et les variétés modulaires que l'on obtient à partir des
structures de Hodge sur $\mathbb{Q}$ sont les
\[
\mathcal{J} / \Gamma , \qquad
\Gamma \subset G(\mathbb{Q}) \cap G(\mathbb{R})^{\circ} \text{ commensurable à } G(\mathbb{Z}) .
\]
Il se peut, note-t-il, qu'il faille imposer aux données des conditions
supplémentaires : arithmétiques, mettant en jeu l'existence d'éléments de
Frobenius dans $I(\mathbb{Q})$ où $I = G/\text{aut. int.}$ ; ou géométriques,
demandant que $G_{\mathbb{R}}$ n'ait pas de facteur compact.

Le cas typique est celui-ci : un $\mathbb{Q}$-espace vectoriel $V$ de dimension
finie muni d'une forme bilinéaire $\varphi$ symétrique ou alternée,
$G = SO(\varphi)$, et $i_{0} : \mathbb{U}_{\mathbb{R}} \to G_{\mathbb{R}}$ avec
$i_{0}(-1) = \mathrm{id}$ si $\varphi$ est symétrique, $-\mathrm{id}$ si elle
est alternée, plus une condition de positivité de $\varphi$ relativement à la
bigraduation de $V_{\mathbb{C}}$ définie par $i_{0}$. Alors $\mathcal{J}$
généralise directement le demi-plan de Siegel, et classe les structures de Hodge
de $V$ compatibles avec $\varphi$ et de nombres de Hodge $\dim
V_{\mathbb{C}}^{p,q}$ donnés. L'idée est que ces $\mathcal{J}$-là sont
« universels » pour tous les autres.

\subsection*{L'argument de Borel, et la contradiction apparente}

Prenons $G(\mathbb{R})$ compact. Alors $\mathcal{J}/\Gamma = \mathcal{J}$ est un
espace homogène projectif complexe $X$ sous $G_{\mathbb{C}}$, par un théorème de
Borel. On obtiendrait ainsi un motif polarisé sur l'espace analytique
$\mathcal{J}$ correspondant à $X$ ; et si l'on admet le yoga GAGA --- venu du cas
des variétés abéliennes polarisées --- selon lequel une telle famille doit aussi
être un motif sur $X$, on obtient une contradiction avec les conjectures de
Tate, puisque $X$ est simplement connexe. D'où la phrase qui ferme le
raisonnement :

\begin{quote}
\emph{So I really do not know what to believe!}
\end{quote}

\subsection*{La condition nécessaire, et la question laissée ouverte}

Revenant au cas général, il ajoute une condition nécessaire. Pour que
$\mathcal{J}/\Gamma$ soit algébrique et que la structure de Hodge $V/\Gamma$
au-dessus le soit aussi, en admettant les conjectures de Tate et de Hodge, il
faut : pour tout $g \in G(\mathbb{R})^{\circ}$ et tout $\Gamma_{0}$ d'indice fini
dans $G(\mathbb{Z})$, le plus petit sous-groupe algébrique $H$ de $G$ tel que
$H_{\mathbb{R}}$ contienne $g\, i_{0}(\mathbb{U}_{\mathbb{R}})\, g^{-1}$
\emph{et} $\Gamma_{0}$ est $G$ tout entier.

Un bloc barré, mais lisible, tire la conséquence : si cette condition était
fausse ne serait-ce que pour l'espace modulaire de Siegel usuel, alors --- sachant
que $\mathcal{J}/\Gamma$ y est algébrique et les $V/\Gamma$ aussi --- il
s'ensuivrait que la conjecture de Hodge ou celle de Tate est fausse.\footnote{Le
bloc est annulé par un quadrillage de traits diagonaux, et le « not » de « I did
not try » est lui-même biffé à l'intérieur du bloc : il a donc, semble-t-il,
essayé. Le résultat de l'essai n'est pas sur le feuillet.}

La lettre se clôt sur deux questions. La première :

\begin{quote}
Si $G_{\mathbb{R}}$ est semi-simple sans facteur compact, $G(\mathbb{Z})$ est-il
toujours Zariski-dense dans $G$ ?
\end{quote}

La réponse est affirmative.\footnote{C'est le théorème de densité de Borel
(1960) : un réseau d'un groupe de Lie semi-simple connexe sans facteur compact
est Zariski-dense ; et $G(\mathbb{Z})$ est un réseau de $G(\mathbb{R})$ par le
théorème de Borel--Harish-Chandra (1962). La page pose la question et ne la
tranche pas ; rien sur le feuillet n'indique qu'il ait fait le rapprochement.}
La seconde porte sur l'exposé de Mumford à Boulder, qu'il dit venir de lire, et
demande si l'image de $\mathcal{J}/\Gamma$ dans la variété modulaire de Siegel
usuelle est une variété algébrique, ou au moins un constructible.

\pagerange{22}{24}
\section*{Connexions intégrables sur un foncteur fibre (pages 22 à 24)}

La liasse suivante s'ouvre sous une chemise de sa main --- « Descriptions
transcendantes de motifs en caractéristique $0$ et de catégories de motifs en
caractéristique $0$ ; doubles réseaux etc. » --- et repose tous ses objets à
neuf. Ces trois feuillets isolent un outil dont les pages 38 à 43 se serviront.

Soit $\mathcal{M}$ une $\otimes$-catégorie sur un corps $k_{0}$, $K$ une
extension de $k_{0}$, et supposons donné un foncteur fibre à valeurs dans
$k_{0}$, de groupe $G$, de sorte que $\mathcal{M} \simeq \operatorname{Rep}(G)$.
On cherche toutes les façons de munir chaque $V \otimes_{k_{0}} K$ d'une
connexion intégrable, fonctorielle en $V$ et compatible au produit tensoriel.

Il en existe déjà une, notée $\Theta_{0}$, déduite des structures dont on
dispose. Toute autre en diffère donc d'un terme :
\[
\Theta_{X} \;=\; \Theta_{0X} + \omega(X) , \qquad X \in \operatorname{Der}(K/k_{0}) .
\]
Autrement dit les connexions forment un espace affine sous l'espace des
$\omega$. Or exiger la compatibilité au produit tensoriel, c'est exiger que
$\omega(X)$ soit une \emph{dérivation du foncteur fibre} ; et une dérivation du
foncteur fibre est donnée par un élément de $\mathfrak{g}_{K} = \mathfrak{g}
\otimes_{k_{0}} K$, où $\mathfrak{g} = \operatorname{Lie} G$. On obtient donc une
application $K$-linéaire
\[
\omega : \operatorname{Der}(K/k_{0}) \longrightarrow \mathfrak{g}_{K} .
\]

L'intégrabilité s'écrit alors, pour $X, Y \in \operatorname{Der}(K/k_{0})$,
\[
\Theta_{0X}\, \omega(Y) - \Theta_{0Y}\, \omega(X) \;=\; \omega([X,Y]) - [\omega(X), \omega(Y)] ,
\]
où $\Theta_{0X}$ et $\Theta_{0Y}$ opèrent sur $\mathfrak{g}_{K}$ par leur action
sur le facteur $K$. C'est l'identité de Maurer--Cartan, et le plus souvent
$\omega$ provient d'un élément de $\Omega^{1}_{K/k_{0}} \otimes_{K}
\mathfrak{g}_{K} = \Omega^{1}_{K/k_{0}} \otimes_{k_{0}} \mathfrak{g}$, voire
déjà d'un $\omega \in \Omega^{1}_{k/k_{0}} \otimes_{k} \mathfrak{g}_{k}$ si $k$
est une extension de type fini de $k_{0}$.

Voici l'usage qu'il veut en faire. Étant donnés $V, W \in \mathcal{M}$ et
$u : V_{K} \to W_{K}$, à quelle condition $u$ provient-il d'un morphisme de
$\mathcal{M}$ ? Il suffit que $u$ commute aux $\Theta_{X}$, donc --- l'autre
terme étant acquis --- aux $\omega(x)$, donc au plus petit sous-espace vectoriel
$\mathfrak{n} \subset \mathfrak{g}$ tel que $\mathfrak{n}_{K}$ contienne les
$\omega(x)$.\footnote{La lettre gothique n'est pas assurée sur le feuillet : le
tracé supporterait aussi bien un $\mathfrak{m}$. Elle ne reparaît pas ailleurs
dans le dossier.} Comme $u$ commute à $\mathfrak{n}$, il suffit donc que
\[
\mathfrak{n} \;=\; \mathfrak{g}
\]
pour conclure. Il note au passage que $\mathfrak{n}$ est déjà une sous-algèbre
de Lie, et qu'on peut la remplacer par l'algèbre de Lie \emph{algébrique}
qu'elle engendre.

C'est, sous une forme tannakienne, la condition que le groupe de Galois
différentiel de la connexion soit $G$ tout entier : la connexion doit être
« assez tordue » pour que rien ne commute avec elle qui ne soit déjà dans
$\mathcal{M}$. Le groupe $R$ engendré par la monodromie, à la page 39, jouera le
même rôle.

\pagerange{26}{31}
\section*{Le torseur des périodes (pages 26 à 31)}

Tout ce qui suit emploie le formalisme tannakien : un groupe se lit sur les
automorphismes d'un foncteur fibre, et une catégorie se reconstitue comme les
représentations de ce groupe. Cela suppose que la catégorie de motifs en jeu
soit tannakienne neutre --- ce que les conjectures standard donneraient, et ce
que le dossier ne fait nulle part. On le dit ici une fois et on ne le répète
plus.\footnote{Le mot « tannakien » n'apparaît nulle part dans le dossier, et
l'axiomatisation qui le porte est postérieure ; le formalisme, lui, y est
employé sans réserve d'un bout à l'autre.}

\subsection*{La donnée}

Soit $k \subset \mathbb{C}$ et $\mathcal{M}$ une catégorie de motifs sur $k$, de
type fini. Elle porte deux foncteurs fibres naturels :
\[
T_{DR} : \mathcal{M} \longrightarrow \operatorname{Mod}_{f}(k) ,
\qquad
T_{B} : \mathcal{M} \longrightarrow \operatorname{Mod}_{f}(\mathbb{Q}) ,
\]
le second passant par les structures de Hodge et dépendant du plongement
$k \subset \mathbb{C}$. Le premier est filtré. Et il existe entre eux un
isomorphisme canonique
\[
\varphi : T_{B} \otimes_{\mathbb{Q}} \mathbb{C} \ \overset{\sim}{\longrightarrow}\ T_{DR} \otimes_{k} \mathbb{C} .
\]

Soit $G = \operatorname{Aut}^{\otimes}(T_{B})$. En comparant $T_{DR}$ à
$T_{B} \otimes_{\mathbb{Q}} k$ on obtient
\[
P \;=\; \operatorname{Isom}^{\otimes}(T_{B} \otimes_{\mathbb{Q}} k,\ T_{DR}) ,
\]
torseur à droite sous $G_{k}$, qui permet de retrouver $T_{DR} = P \times^{G_{k}}
(T_{B} \otimes_{\mathbb{Q}} k)$ ; et l'isomorphisme $\varphi$ définit un élément
canonique
\[
\varphi \in P(\mathbb{C}) .
\]
C'est la matrice des périodes, et c'est le point transcendant autour duquel
tourne tout le dossier. La connaissance de $G$, de $P$ et de $\varphi$ permet de
retrouver $\mathcal{M}$ avec ses deux foncteurs fibres et leur comparaison.

\subsection*{Le filtré, le gradué, et les deux cocaractères}

Pour retrouver la filtration de $T_{DR}$ et en déduire la structure de Hodge sur
$T_{B}$, il faut faire entrer le gradué. Posons $T_{\mathrm{Hdg}} =
\operatorname{gr}^{*} T_{DR}$, et
\[
G' \;=\; P \times^{G_{k}} G_{k} \;=\; \operatorname{ad}(P) \;\simeq\; \operatorname{Aut}^{\otimes}(T_{DR}) .
\]
Soit $H' \subset G'$ le groupe des automorphismes de $T_{DR}$ qui respectent la
filtration et induisent l'identité sur le gradué : il est unipotent. Soit
\[
Q \;=\; \operatorname{Isom}^{\otimes}_{\ \mathrm{id}\ \mathrm{sur}\ \operatorname{gr}^{*}}(T_{DR},\ T_{\mathrm{Hdg}}) ,
\]
torseur à droite sous $H'$, et posons $H'' = \operatorname{ad}(H')$ et
\[
G'' \;=\; Q \times^{H'} G' \;\simeq\; \operatorname{Aut}^{\otimes}(T_{\mathrm{Hdg}}) .
\]
La bigraduation de $T_{\mathrm{Hdg}}$ définit alors deux cocaractères
\[
i_{1}, i_{2} : \mathbb{G}_{m,k} \longrightarrow G''
\]
soumis à
\[
\varepsilon\, i_{1}(\lambda) = \varepsilon\, i_{2}(\lambda) = \lambda ,
\qquad i_{1}(\lambda)\, i_{2}(\lambda) = i(\lambda) ,
\]
les deux cocaractères étant les deux degrés $p$ et $q$ et leur produit le degré
total.\footnote{Le feuillet 27 écrit le but de ces deux morphismes $G'$ ; c'est
$G''$ qu'il vient d'identifier à $\operatorname{Aut}^{\otimes}(T_{\mathrm{Hdg}})$
deux lignes plus haut, et c'est $G''$ qui est écrit ici.}

\subsection*{Les données $a)$ à $g)$, et ce qu'on ne sait pas en dire}

Le tout se résume en une liste :

\begin{itemize}
  \item[$a)$] un groupe algébrique $G$ sur $\mathbb{Q}$, avec
  $\mathbb{G}_{m} \xrightarrow{\ i\ } G \xrightarrow{\ \varepsilon\ } \mathbb{G}_{m}$ ;
  \item[$b)$] un groupe algébrique $G''$ sur $k$, avec
  $\mathbb{G}_{m} \xrightarrow{\ i''\ } G'' \xrightarrow{\ \varepsilon''\ } \mathbb{G}_{m}$ ;
  \item[$c)$] les cocaractères $i_{1}, i_{2}$ soumis aux relations ci-dessus ;
  \item[$d)$] un torseur $Q$ sous le groupe unipotent $H'' \subset G''$ ;
  \item[$e)$] un torseur $P$ sous $G'$ ;
  \item[$f)$] un isomorphisme $G \otimes_{\mathbb{Q}} k \simeq \operatorname{ad}_{G'}(P)$ ;
  \item[$g)$] un élément $\varphi \in P(\mathbb{C})$, compatible aux structures
  $i$ et $\varepsilon$.
\end{itemize}

Inversement, de ces données on récupère $\mathcal{M}$ comme
$\operatorname{Rep}(G)$, le foncteur $T_{B}$ avec sa structure de Hodge,
$T_{DR}$ comme $P \times^{G} (T_{B} \otimes_{\mathbb{Q}} k)$, le gradué
$T_{\mathrm{Hdg}} \simeq T_{DR} \times^{H'} Q$, et la bigraduation par $i_{1}$,
$i_{2}$ ; $\varphi$ transporte enfin la filtration de $T_{DR} \otimes_{k}
\mathbb{C}$ sur $T_{B} \otimes_{\mathbb{Q}} \mathbb{C}$.

Trois conditions sont à imposer : $1)$ que pour toute représentation fidèle $V$
la filtration obtenue sur $V_{\mathbb{C}}$ en fasse un modèle de Hodge ; $2)$ que
ce dernier soit \emph{intrinsèque} ; $3)$ que la restriction des scalaires de
$\mathbb{C}$ à $k$ déduite de $\varphi$ corresponde à une restriction des
scalaires du motif. Et il écrit aussitôt ce qui est le vrai état des lieux :

\begin{quote}
On ne sait pas expliciter $2)$ et $3)$ directement en termes des données
$a)$ à $g)$.
\end{quote}

La seule condition $1)$ permet en revanche de définir un élément
$\psi \in Q(\mathbb{C})$, par la condition que la bigraduation correspondante
satisfasse $\overline{\psi} = \psi^{-1}$ --- c'est-à-dire la symétrie
$V^{p,q} = \overline{V^{q,p}}$, transportée dans le torseur.

\subsection*{Le facteur $2i\pi$}

Si les données $a)$ à $g)$ proviennent effectivement d'une catégorie de motifs
sur $k$, il y a dans $P \times^{G} \mathbb{G}_{m}$ un isomorphisme marqué
$\alpha$, et, considérant l'élément $\dot{\varphi} \in (P \times^{G}
\mathbb{G}_{m})(\mathbb{C})$ induit par $\varphi$, on doit avoir
\[
\dot{\varphi} \;=\; \frac{1}{2 i \pi}\, \alpha_{\mathbb{C}} .
\]
C'est la normalisation du motif de Tate : la composante de la matrice des
périodes sur l'objet inversible $\mathbb{Q}(1)$ vaut $2i\pi$, et pas autre chose.
Et c'est, dit-il, une \emph{restriction arithmétique} sur les données $a)$ à
$g)$ : toutes les configurations abstraites de groupe, torseur et point complexe
ne proviennent pas de motifs, et la première obstruction est déjà visible sur le
plus petit motif qui soit.

\subsection*{La question inverse}

Partons cette fois d'une structure de Hodge $(V_{\mathbb{Q}}, V_{\mathbb{C}} =
\bigoplus V^{p,q})$, soit $G$ son groupe de Mumford--Tate,\footnote{Le nom est
de sa main, au feuillet 30.} et soit $V_{k} \subset V_{\mathbb{C}}$ une
$k$-structure compatible avec la filtration naturelle. Quand ce triplet
provient-il d'un motif sur $k$ ?

Il faut, dit-il, que $G'_{\mathbb{C}}$ soit définissable sur $k$ via
$V_{k} \subset V_{\mathbb{C}}$ ; et lorsque $k$ est algébriquement clos, que
$G'_{\mathbb{C}} = G_{\mathbb{C}}$. Concrètement : tout élément de
$\bigotimes_{\mathbb{Q}}(V^{\vee} \otimes_{\mathbb{Q}} V)$ invariant sous $G$ et
de bidegré $0$ doit se retrouver dans $\bigotimes_{k}(V^{\vee}_{k} \otimes_{k}
V_{k})$. Ce sont des conditions très restrictives, et il ajoute qu'elles ne sont
certainement pas suffisantes : pour $k = \mathbb{Q}$ et $V_{k} = V_{\mathbb{Q}}$
ce n'est en général pas possible. D'où la phrase du feuillet 31, qui est tout ce
qu'il en dit :

\begin{quote}
Il faut donc compléter les conditions nécessaires précédentes par des
conditions du type « argument de transcendance ».
\end{quote}

C'est exactement l'endroit où l'on attend aujourd'hui la conjecture des
périodes : les conditions algébriques ne suffisent pas, et ce qui manque est une
minoration du degré de transcendance des périodes.\footnote{Le dossier ne
formule aucune telle conjecture ; il en marque la place et passe.}

\pagerange{33}{34}
\section*{La situation abstraite des doubles réseaux (pages 33 et 34)}

Ces deux feuillets démontrent, dans l'abstrait, le critère dont les pages 45 à
49 auront besoin. La question est celle-ci : deux foncteurs fibres, l'un sur
$k$, l'autre sur une extension $K$, deviennent isomorphes après extension à un
corps $L$ contenant les deux ; que reste-t-il, dans un objet, de commun aux deux
structures rationnelles ?

Soient donc $\mathcal{M}$ une $\otimes$-catégorie sur $k$, $T$ un foncteur fibre
sur $k$, $T'$ un foncteur fibre sur $K$, et $\varphi : T \otimes_{k} L \simeq
T' \otimes_{K} L$ un $\otimes$-isomorphisme. Posons
\[
G = \operatorname{Aut}^{\otimes}(T) , \qquad
P = \operatorname{Isom}^{\otimes}(T_{K}, T') , \qquad
\varphi \in P(L) ,
\]
$P$ étant un torseur à droite sous $G_{K}$. Ces données reconstituent
$\mathcal{M}$, $T$, $T'$ et $\varphi$.

Supposons d'abord $P$ trivial, et soit $\xi \in P(K)$. On a $\varphi = \xi g$
avec $g \in G(L)$, et changer $\xi$ en $\xi a$ change $g$ en $ga$ : ce qui est
canonique, c'est donc
\[
\dot{g} \in G(L)/G(K) ,
\]
et les couples $(G, \dot{g})$ équivalent aux triplets $(G, P, \varphi)$. Pour un
$G$-module $V$, le réseau $T'(V) \subset V \otimes_{k} L$ n'est autre que
$g\,(V \otimes_{k} K)$, de sorte que
\[
V \cap T'(V) \;=\; \{\, x \in V \ \mid\ g^{-1} x \in V \otimes_{k} K \,\} .
\]

Soit $T_{K} \subset G_{K}$ l'adhérence de l'orbite de $g$, et soit $R$
l'adhérence de l'image de
\[
T_{K} \times T_{K} \longrightarrow G_{K} , \qquad (s,t) \longmapsto t s^{-1} ,
\]
qui ne dépend pas du choix de $g$ dans sa classe. On a alors, pour $x \in V$,
\[
x \in V \cap T'(V) \quad \Longleftrightarrow \quad R \subset G_{x,K} ,
\]
où $G_{x}$ est le stabilisateur de $x$. Soit enfin $S$ l'image de $R$ dans $G$ et
$H$ le sous-groupe algébrique de $G$ engendré par $S$ ; comme $T_{K}$ est
fidèlement plat sur $K$, la descente donne
\[
V \cap T'(V) \;=\; V^{H} .
\]

Le cas d'un $P$ quelconque se ramène au précédent sans point-base : le morphisme
de bitorseurs $P \times_{K} P \to G_{K}$, $(s,u) \mapsto s^{-1}u$, a une image
fermée $R$ définie intrinsèquement, et la même descente fidèlement plate donne
la même égalité.\footnote{Le feuillet écrit ici $V \cap T'(W) = V^{H}$ ; c'est
$T'(V)$ que l'argument porte, et c'est $V$ qui est écrit ici. Le $W$ est bien sur
le feuillet.}

Comme $H \subset G$, on a toujours $V^{G} \subset V^{H}$, et l'on obtient donc le
critère, qui est ce que les deux feuillets cherchaient :

\begin{quote}
Le foncteur $\operatorname{Rep}(G) = \mathcal{M} \to \{\text{doubles réseaux}\}$
est pleinement fidèle si et seulement si, pour tout $V$,
\[
V^{H} \;=\; V^{G} .
\]
\end{quote}

Autrement dit : la pleine fidélité échoue exactement dans la mesure où le groupe
engendré par la position relative des deux réseaux est plus petit que le groupe
de Galois tout entier. C'est la forme abstraite de la question que la page 47
posera des motifs.

\pagerange{36}{37}
\section*{La récapitulation adélique (pages 36 et 37)}

Ce feuillet est une page de tableau : sept encadrés étiquetés $A)$ à $G)$, puis
neuf propriétés numérotées. On le donne dans son ordre.

Soit $K \subset \mathbb{C}$, $\overline{K}$ la clôture algébrique de $K$ dans
$\mathbb{C}$, $\pi_{1} = \operatorname{Gal}(\overline{K}/K)$ et $\mathbb{A}$
l'anneau des adèles finis de $\mathbb{Q}$. On considère une catégorie de motifs
de type fini sur $K$ et le groupe $G$ associé au foncteur de Betti. Les données
sont :

\begin{itemize}
  \item[$A)$] le groupe $G$ ;
  \item[$B)$] la chaîne $G''_{K} \supset H''_{K} \cdots Q \cdots H'_{K} \subset
  G'_{K} \cdots P \cdots G_{K}$, avec $i_{1}, i_{2} : \mathbb{G}_{m}
  \rightrightarrows G''_{K}$ --- c'est-à-dire toute la donnée des pages 26 et 27 ;
  \item[$C)$] un morphisme $u : \pi_{1} \to G(\mathbb{A})$ ;
  \item[$D)$] une connexion absolue intégrable invariante sous $G_{K}$ ---
  celle des pages 22 à 24 ;
  \item[$E)$] une donnée de « densité », soumise à la compatibilité avec $A)$,
  $B)$, $C)$ ;
  \item[$F)$] la commutativité du carré reliant $u$ et $v$ ;
  \item[$G)$] un morphisme $w : \pi_{1}(S(\mathbb{C}), \xi) \to G(\mathbb{C})$.
\end{itemize}

Du côté complexe on a la connexion absolue $\omega \in
\Omega^{1}_{\mathbb{C}/\mathbb{Q}} \otimes_{\mathbb{Q}} \mathfrak{g}$, les
isomorphismes $G' \otimes_{K} \mathbb{C} \simeq G_{\mathbb{C}}$ et
$G'' \otimes_{K} \mathbb{C} \simeq G_{\mathbb{C}}$, d'où deux morphismes
$i^{\mathbb{C}}_{1}, i^{\mathbb{C}}_{2} : \mathbb{G}_{m,\mathbb{C}} \to
G_{\mathbb{C}}$ ; puis, pour un modèle rigidifié convenable $S$ de $K$, de type
fini sur $\mathbb{Q}$, le morphisme $w : \pi_{1}(S(\mathbb{C}), \xi) \to
G(\mathbb{C})$ ; enfin $P/G^{0}$, torseur sur $K$ de groupe $G/G^{0}$, ponctué
par $\xi : K \hookrightarrow \mathbb{C}$, d'où
$v : \pi_{1} \to (G/G^{0})(\mathbb{Q})$.

Les deux carrés dont il affirme la commutativité sont ceux-ci.

\begin{tikzcd}[column sep=small, row sep=small, nodes={font=\scriptsize}]
\pi_{1}(S(\mathbb{C}), \xi) \arrow[r, "w"] \arrow[d, "\mathrm{can.}"'] & G(\mathbb{Q}) \arrow[d, "\mathrm{can.}"] \\
\pi_{1}(S, \xi) \arrow[r, "u"] & G(\mathbb{A})
\end{tikzcd}

\begin{tikzcd}[column sep=small, row sep=small, nodes={font=\scriptsize}]
\pi_{1} \arrow[r, "v"] \arrow[d, "u"'] & (G/G^{0})(\mathbb{Q}) \arrow[d] \\
G(\mathbb{A}) \arrow[r] & (G/G^{0})(\mathbb{A})
\end{tikzcd}

Les propriétés, telles qu'il les numérote :

\begin{enumerate}
  \item[$1)$] $i^{\mathbb{C}}_{2} = i^{\mathbb{C}}_{1}$, de sorte que $\alpha$
  est connu dès qu'on connaît l'un des deux ;
  \item[$2)$] $\pi_{1}(S(\mathbb{C}), \xi) \to G(\mathbb{C})$ se factorise par
  $G(\mathbb{Q})$ ;
  \item[$3)$] $\pi_{1} \to G(\mathbb{A})$ se factorise par
  $\pi_{1}(S, \xi) \to G(\mathbb{A})$ pour $S$ convenable ;
  \item[$4)$] les deux carrés ci-dessus sont commutatifs ;
  \item[$5)$] le composé $\pi_{1}(S, \xi) \to G(\mathbb{A})
  \xrightarrow{\ \varepsilon\ } \mathbb{A}^{*}$ est le morphisme donné par
  l'action sur les racines de l'unité ;
  \item[$6)$] $v$ est un épimorphisme ;
  \item[$7)$] pour tout point fermé $s$ de $S$, l'élément de Frobenius $f_{s}$
  a une image dans $G(\mathbb{A})$ --- « redonne Weil » ;
  \item[$8)$] $G_{\mathbb{R}}$, muni de $i^{\mathbb{C}}_{1}$ et
  $i^{\mathbb{C}}_{2}$, est un groupe de Hodge ;
  \item[$9)$] pour tout $\ell$, l'image de $\pi_{1}$ dans $G(\mathbb{Q}_{\ell})$
  est un sous-groupe ouvert, et $V \otimes_{\mathbb{Q}} \mathbb{Q}_{\ell}$ est
  \emph{semi-simple}.
\end{enumerate}

Cette liste est le programme entier de la théorie de Galois motivique sur une
base, écrit sur un feuillet. La propriété $5)$ est le caractère cyclotomique ;
$7)$ est ce qui doit redonner les conjectures de Weil ; $8)$ est le pont avec
les pages 3 à 11, et dit que la théorie de Hodge du début du dossier est la
composante archimédienne de ce tableau ; $9)$ est l'ouverture de l'image
$\ell$-adique jointe à la semi-simplicité. Il ferme sur deux mots :
« devra être précisé ».\footnote{Le feuillet 37 ne porte que deux lignes de
notation isolées --- $G/Q$, puis $G'_{K}\ P\ G_{K}\ G_{\mathbb{C}}/\mathbb{C}$
--- sans phrase autour ; ce sont les objets du tableau, rien de plus.}

\pagerange{38}{43}
\section*{La donnée transcendante et le revêtement équivariant (pages 38 à 43)}

\subsection*{L'application des périodes}

Il introduit alors ce qu'il appelle une structure « de nature transcendante ».
Soit $\overline{S}_{\mathbb{C}}$ le revêtement universel de $S(\mathbb{C})$ en
$\xi$, et soit $\mathcal{J}$ l'ensemble des morphismes $\mathbb{S} \to
G_{\mathbb{R}}$ correspondant au morphisme central donné
$i(\lambda) = i_{1}(\lambda) i_{2}(\lambda)$ --- c'est-à-dire, exactement,
l'espace des structures de Hodge sur $G_{\mathbb{R}}$ construit à la page 7 et
retrouvé à la page 15.\footnote{Il le note ici $\Gamma_{G}$, et $\mathcal{J}$ aux
pages 7 à 19 ; c'est un seul objet et il porte ici un seul nom. Rien sur les
feuillets ne signale l'identité, qui est pourtant ce qui fait du dossier un seul
texte.} On a alors une application
\[
f : \overline{S}_{\mathbb{C}} \longrightarrow \mathcal{J}
\]
soumise à trois conditions :

\begin{enumerate}
  \item[$a)$] $f(\overline{\xi})$ est le morphisme $\mathbb{S} \to
  G_{\mathbb{R}}$ défini par $i^{\mathbb{C}}_{1}$ et $i^{\mathbb{C}}_{2}$ ;
  \item[$b)$] $f(g \overline{s}) = w(g)\, f(\overline{s})$ pour
  $g \in \pi_{1}(S(\mathbb{C}), \xi)$ ;
  \item[$c)$] $f$ est \emph{holomorphe}.
\end{enumerate}

C'est l'application des périodes d'une variation de structures de Hodge : $b)$
est l'équivariance sous la monodromie, $c)$ l'holomorphie, et la compatibilité
au torseur $P$ énoncée plus bas tient lieu de la condition de transversalité.

Vient ensuite, moyennant la conjecture de Hodge, un énoncé de rigidité : le
composé $\overline{S}_{\mathbb{C}} \to \mathcal{J} \to \mathcal{J}_{\widetilde{G}}$
est une application \emph{constante}, où $\widetilde{G}$ est le quotient
introduit ci-dessous.\footnote{Le feuillet 38 numérote $13)$ deux énoncés
consécutifs ; ils sont distingués ici, la discordance étant sur la page. Le
second dit que, pour $\overline{s}$ au-dessus d'un point $t$ de $S$, le morphisme
$f(\overline{s})_{\mathbb{C}}$ --- qui définit une bigraduation, donc une
filtration, sur les $G$-modules --- est compatible avec le torseur $P_{t}$ sous
$G_{k(t)}$ et sa $\mathbb{C}$-trivialisation.}

\subsection*{Le groupe de monodromie, et la descente à $k$}

Soit $R \subset G$ le sous-groupe algébrique engendré par
$w(\pi_{1}(S(\mathbb{C}), \xi))$ : c'est le groupe de monodromie. Il est
\emph{invariant} dans $G$, et c'est cette normalité qui rend le quotient
utilisable.\footnote{On reconnaît ici l'énoncé qui deviendra le « théorème de la
partie fixe » et la normalité du groupe de monodromie connexe dans le groupe de
Mumford--Tate générique. Le feuillet en fait dépendre la démonstration de la
propriété $9)$ du tableau, c'est-à-dire de la semi-simplicité $\ell$-adique ;
c'est une implication du dossier et elle n'est pas établie.} Posant
$\widetilde{G} = G/R$, on obtient
\[
\widetilde{u} : \pi_{1}(\operatorname{Spec} k, \xi) \longrightarrow \widetilde{G}(\mathbb{A}) ,
\]
et l'on considère le torseur $\widetilde{P} = P/R$ sous $\widetilde{G}$, muni de
la connexion intégrable déduite de celle de $P$. Cette connexion se descend à
$k$ ; donc, moyennant Hodge, le couple $(\widetilde{P}, \widetilde{G})$ avec sa
connexion provient d'un couple $(\widetilde{P}_{0}, \widetilde{G}_{k})$ sur $k$,
avec un isomorphisme de torseurs
\[
\widetilde{P} \;=\; P/R \ \simeq\ \widetilde{P}_{0} \otimes_{k} K .
\]
Il en va de même de $\widetilde{G}' = G'/R'$, de $\widetilde{H}'$, du torseur
$\widetilde{Q}$ et des cocaractères $\widetilde{i}_{1}, \widetilde{i}_{2}$, qui
descendent tous en des objets sur $k$ : la donnée horizontale pour la connexion
est exactement la donnée de descente.

\subsection*{Sur une base : ce que la donnée d'un point réduit}

La page 40 pose le cadre définitif. Soit $S$ un schéma régulier de type fini sur
$\mathbb{Q}$, $k$ la clôture algébrique de $\mathbb{Q}$ dans
$\Gamma(S, \mathcal{O}_{S})$ --- de sorte que $S$ est géométriquement connexe
sur $k$ --- et $k \hookrightarrow \mathbb{C}$ un plongement. Soit $\mathcal{M}$
une catégorie de motifs lisses de type fini sur $S$. Chaque point
$\xi \in S(\mathbb{C})$ donne un foncteur fibre $T_{\xi} : M \mapsto
T_{B}(M_{\xi})$ ; posant $G = \operatorname{Aut}^{\otimes}(T)$, on a
$\mathcal{M} \simeq \operatorname{Rep}(G)$, et
\[
P = \operatorname{Isom}^{\otimes}(T_{S},\ T_{DR})
\]
est un torseur sous $G_{S}$, muni d'une connexion intégrable relative à $k$ et
d'une $\otimes$-filtration sur $T_{DR}$.

Fixons maintenant $\xi_{0} \in S(\mathbb{C})$ et prenons $T = T_{\xi_{0}}$.
Alors la donnée des revêtements équivariants équivaut à celle de
\[
w : \pi_{1}(S(\mathbb{C}), \xi_{0}) \longrightarrow G(\mathbb{Q}) \subset G(\mathbb{C})
\]
et d'un point $x_{0}$ de $P_{\xi_{0}}$, c'est-à-dire d'une trivialisation de
$P_{\xi_{0}}$. Ce point définit un unique morphisme
$q : \overline{S}_{\mathbb{C}} \to P(\mathbb{C})$ envoyant le point marqué sur
$x_{0}$, soumis à
\[
q(\overline{s}\,\sigma) = q(\overline{s}) \cdot w(\sigma) ,
\]
et l'on obtient
\[
\widetilde{P}_{\mathbb{C}} \;=\; \overline{S}_{\mathbb{C}} \times^{\Gamma} G(\mathbb{Q})
\ \hookrightarrow\ P(\mathbb{C}) , \qquad \Gamma = \pi_{1}(S(\mathbb{C}), \xi_{0}) ,
\]
qui est la restriction du groupe structural de $G(\mathbb{C})$ à
$G(\mathbb{Q})$. L'application des périodes se relit sur
$\widetilde{P}_{\mathbb{C}}$ :
\[
f : \widetilde{P}_{\mathbb{C}} \longrightarrow \mathcal{J} ,
\qquad f(x \cdot g) = g^{-1} f(x) \ \ (g \in G(\mathbb{Q})) ,
\]
holomorphe, et compatible en chaque point $s$ au torseur $P_{s}$ sous
$G_{k(s)}$ et à sa trivialisation. Se donner enfin la couche adélique, c'est se
donner une factorisation de $\Gamma \to G(\mathbb{Q}) \to G(\mathbb{A})$ par le
groupe fondamental algébrique $\pi_{1}(S, \xi_{0})$.

\subsection*{Le bilan, et les deux questions}

La page 43 résume le tout en cinq données :

\begin{itemize}
  \item[$A)$] un groupe réductif $G$ sur $\mathbb{Q}$ ;
  \item[$B)$] un fibré principal à connexion intégrable $P$ sous $G_{S}$ ;
  \item[$C)$] la restriction du groupe structural à $G(\mathbb{Q})$, d'où
  $\widetilde{P} \hookrightarrow P(\mathbb{C})$ ;
  \item[$D)$] l'application holomorphe équivariante $f : \widetilde{P} \to
  \mathcal{J}$, compatible avec $B)$ et $C)$ ;
  \item[$E)$] la restriction de $\widetilde{P} \times^{G(\mathbb{Q})}
  G(\mathbb{A})$ à $S$, compatible verticalement à $G/G^{0}$.
\end{itemize}

Et il note que la connaissance de $A)$, $B)$, $C)$ et du seul point
$f(\xi_{0})$ --- au-dessus des points génériques de $S$ --- doit en principe
permettre de reconstituer $D)$ tout entier, et les données venant de $E)$.
D'où les deux questions sur lesquelles la section s'arrête :

\begin{quote}
Formule analytique de prolongement de $f$ ? \\
Relations entre $D)$ et $E)$ ?
\end{quote}

La seconde porte, en marge, une glose entre crochets : « Jugendtraum
? ».\footnote{La lecture du mot est donnée sous réserve par la transcription. Si
elle est juste, elle désigne le rêve de jeunesse de Kronecker --- la description
explicite des extensions abéliennes par des valeurs de fonctions transcendantes
--- et la question est alors celle de la loi de réciprocité liant la donnée
transcendante $D)$ à la donnée adélique $E)$ : c'est-à-dire, en termes
d'aujourd'hui, la théorie des modèles canoniques. Le reste de la marge n'est pas
lisible.}

\pagerange{45}{49}
\section*{Foncteurs pleinement fidèles (pages 45 à 49)}

Ces cinq feuillets font la liste des foncteurs de réalisation candidats et
demandent de chacun s'il est pleinement fidèle, et à quel prix. C'est la partie
du dossier la plus près d'un programme.

Soit $S$ un $\mathbb{Q}$-schéma régulier connexe muni d'un point géométrique
$\xi : \operatorname{Spec}(\mathbb{C}) \to S$. Le foncteur
$\xi^{*} : \operatorname{Mot}(S) \to \operatorname{Mot}(\mathbb{C})$ est fidèle,
et l'on regarde le composé

\begin{tikzcd}[column sep=small, row sep=small, nodes={font=\scriptsize}]
\operatorname{Mot}(S) \arrow[r, "\xi^{*}"] \arrow[rrr, bend right=25, "(T_{B})_{\xi}"'] & \operatorname{Mot}(\mathbb{C}) \arrow[r, "T_{B+H}"] & \operatorname{Hodge}(\mathbb{Q}, \mathbb{C}) \arrow[r] & \operatorname{Mod}_{f}(\mathbb{Q})
\end{tikzcd}

\subsection*{Par Hodge}

$1)$ Soit $\xi$ localisé au point générique $y$ de $S$. À tout motif $M$ sur $S$
on associe le triplet
\[
\{\, T_{DR}(M_{y}),\ T_{B}(M_{\xi}),\ \varphi_{M} \,\} ,
\qquad
\varphi_{M} : T_{B}(M_{\xi}) \otimes_{\mathbb{Q}} \mathbb{C} \ \simeq \ T_{DR}(M_{y}) \otimes_{k(y)} \mathbb{C} ,
\]
en considérant $T_{DR}(M_{y})$ comme \emph{filtré}. Le foncteur ainsi obtenu est
pleinement fidèle, moyennant la conjecture de Hodge simple. On a le même
résultat, $1\,b)$, en remplaçant le filtré par le bigradué
$T_{\mathrm{Hdg}}(M)$ et $\varphi_{M}$ par l'isomorphisme $\psi_{M}$
correspondant.

\subsection*{Par Tate}

$2)$ Soit $\pi = \pi_{1}(S, \xi)$. À tout $M$ on associe
\[
\{\, T_{B}(M_{\xi}),\ \rho_{\ell, M} \,\} ,
\qquad
\rho_{\ell, M} : \pi \longrightarrow \operatorname{Aut}_{\mathbb{Q}_{\ell}}\bigl( T_{B}(M_{\xi}) \otimes_{\mathbb{Q}} \mathbb{Q}_{\ell} \bigr) ,
\]
c'est-à-dire le réseau de Betti et la représentation galoisienne, \emph{sans}
structure de Hodge.\footnote{Un « $+H$ » écrit après le $T_{B}$ a été biffé sur
le feuillet : il écarte ici délibérément la structure de Hodge, qu'il
rétablira quelques lignes plus bas sous le nom $T_{B-H}$.} Le foncteur est
pleinement fidèle si $S$ est de type fini sur $\mathbb{Q}$, en admettant les
conjectures de Tate et en utilisant un résultat de Griffiths.

Si $S$ est localement de type fini sur $\mathbb{C}$, on peut aussi prendre
$\pi' = \pi_{1}(S(\mathbb{C}), \xi)$ et la représentation $\rho_{M}$ à valeurs
dans $\operatorname{Aut}_{\mathbb{Q}}$, à condition de remplacer $T_{B}$ par
$T_{B-H}$, c'est-à-dire de rétablir la structure de Hodge --- et c'est alors la
conjecture de Hodge qu'on utilise, non celle de Tate. Le couple de variantes est
instructif : la même pleine fidélité s'achète, selon le corps de base, avec l'une
ou l'autre des deux conjectures.

$3)$ À tout $M$ on associe
\[
\bigl( T_{DR}(M),\ \sigma,\ T_{B}(M_{\xi}),\ \varphi \bigr) ,
\]
où $T_{DR}(M)$ est la cohomologie de De Rham de $M$ sur $S$ vue comme espace
\emph{filtré}, $\sigma$ sa stratification de Gauss--Manin relativement à un corps
de base algébriquement clos $k = \overline{\mathbb{Q}}$, et $\varphi$
l'isomorphisme canonique. Le foncteur est pleinement fidèle moyennant Hodge, en
utilisant Griffiths.

\subsection*{Le « double niveau »}

$4)$ Supposons $S = \operatorname{Spec} k$ avec $k$ extension
\emph{algébrique} de $\mathbb{Q}$, et considérons le foncteur
\[
(4.1) \qquad M \longmapsto \{\, T_{DR}(M),\ T_{B}(M),\ \varphi \,\} ,
\]
où l'on ne considère \emph{ni} $T_{DR}(M)$ comme filtré, \emph{ni} $T_{B}(M)$
comme muni d'une structure de Hodge. Ce foncteur est-il pleinement fidèle ?
C'est la question centrale de la section, et elle est posée sans réponse.

Elle se ramène aussitôt au cas $k = \overline{\mathbb{Q}}$, et devient alors
équivalente à la pleine fidélité de
\[
(4.2) \qquad M \longmapsto \{\, T_{B}(M),\ \sigma \,\} ,
\]
$\sigma$ étant la stratification canonique sur
$T_{B}(M) \otimes_{\mathbb{Q}} \mathbb{C}$ --- car $\sigma$ et le réseau
$T_{DR}(M) \subset T_{B}(M) \otimes \mathbb{C}$ se déterminent mutuellement,
$T_{DR}(M)$ étant l'ensemble des éléments horizontaux pour $\sigma$.

Voici la traduction dans le groupe. Soit
$G = \underline{\operatorname{Aut}}^{\otimes}(T_{B})$ le groupe de Galois de
$\operatorname{Mot}(k)$. La donnée revient à celle d'un
\[
\dot{g} \in G(\mathbb{C})/G(\overline{\mathbb{Q}}) ,
\]
et si $V$ est le $G$-module correspondant à $M$, alors $T_{DR}(M)$ n'est autre
que le $\overline{\mathbb{Q}}$-réseau $g(V_{\overline{\mathbb{Q}}}) \subset
V \otimes_{\mathbb{Q}} \mathbb{C}$. La question devient donc simplement : pour
tout $V$, l'inclusion
\[
V^{G} \ \subset \ V \cap g(\overline{V})
\]
est-elle une \emph{égalité} ?\footnote{Le sens de l'inclusion est celui du
feuillet et il compte : $V^{G} \subset V \cap g(\overline{V})$ vaut toujours, un
invariant de tout le groupe de Galois étant à la fois dans le réseau de Betti et
dans celui de De Rham. C'est l'égalité qui est la question, et l'écrire comme une
identité serait affirmer ce qui est demandé.} On peut se ramener à $V$ simple ;
et si $V$ est de dimension $1$, c'est une puissance de Tate et cela marche.

Le feuillet 48 pousse la réduction : la relation $x \in \overline{V} \cap
g(\overline{V})$ équivaut à $\overline{R} \subset G_{x}$, où $G_{x}$ est le
stabilisateur de $x$ et $\overline{R}$ l'adhérence de l'image de
$\overline{T} \times \overline{T} \to \overline{G}$, $(s,t) \mapsto s^{-1}t$. On
retrouve exactement le groupe $H$ des pages 33 et 34, et la question est celle,
déjà posée là, de savoir si $V^{H} = V^{G}$.

$5)$ Sans l'hypothèse que $k$ est algébrique sur $\mathbb{Q}$, le foncteur
$(4.1)$ n'est plus pleinement fidèle en général : si $k$ est de type fini sur
$\mathbb{Q}$ et $M$ fixé, $T_{DR}(M)$ finit par contenir $T_{B}(M)$, ce dont il
écrit en un mot : « idiot ! ». Tout au plus peut-on espérer, à $M$ fixé, un
résultat de disproportion entre $T_{B}$ et $T_{DR}$ si $k$ est modéré ---
par exemple en termes de degré de transcendance minimum.

$6)$ En revanche, sans hypothèse sur $k$ sauf d'être algébriquement clos, il
semble encore raisonnable d'espérer que $(4.2)$ soit pleinement fidèle, ce qui
ramène au cas $k = \mathbb{C}$. Et le foncteur de $3)$ \emph{privé de la
filtration} sur $T_{DR}(M)$ serait pleinement fidèle si la conjecture du double
niveau l'était.

La section laisse donc une conjecture, sous le nom qu'il lui donne --- le
« double niveau » --- et un critère exact pour elle, $V^{H} = V^{G}$, démontré
onze feuillets plus tôt dans l'abstrait. Ce qu'elle ne fait pas, c'est décider.

\pagerange{50}{51}
\section*{Les classes de De Rham essentiellement algébriques (pages 50 et 51)}

Deux feuillets, et le dossier s'arrête.

Soit $K \subset \mathbb{C}$ un sous-corps et $M$ un motif sur $K$, d'où
$T_{DR}(M) = \operatorname{Hom}(1, M) \otimes_{\mathbb{Q}} K$. Un sous-module de
$T_{DR}(M)$ provenant des classes algébriques satisfait deux conditions
nécessaires :

\begin{enumerate}
  \item[$a)$] être contenu dans la partie de poids $0$,
  $T_{DR}(M)^{00}_{\mathbb{C}}$ ;
  \item[$b)$] être contenu dans $\bigl( T_{DR}(M)_{\mathbb{C}} \bigr)_{\mathbb{R}}$.
\end{enumerate}

A-t-on une réciproque ? Il examine trois cas.

$1^{\circ}$ Pas pour $a)$ et $b)$ simultanément. Si $G$ est de poids $0$ la
condition $a)$ est vide, donc sa réciproque n'est vraie que si $G$ est trivial ;
et un exemple galoisien de groupe $G$ sur $\mathbb{Q}$ rend $b)$ vide à son tour.
Il faut donc supposer $G$ connexe, et l'on peut se ramener à
$K = \overline{\mathbb{Q}}$.

$2^{\circ}$ Pour la réciproque de $a)$ seule, il faut un groupe de Galois
motivique de dimension $\geqslant 3$. Fixant le plongement
$K \hookrightarrow \mathbb{C}$, les deux cocaractères $j_{1}, j_{2} : G_{K}
\rightrightarrows G'$ définissent la bigraduation, et une représentation fidèle
$V'$ donnera
\[
{V'}^{\,T} \;=\; T_{DR} \cap (T_{DR})^{00}_{\mathbb{C}} \ \neq\ {V'}^{\,G'} ,
\]
où $T$ est l'image de $(j_{1}, j_{2}) : G_{K} \times G_{K} \to G'$ : c'est-à-dire
un contre-exemple, de la forme même du critère des pages 33 et 34, un groupe
strictement plus petit ayant strictement plus d'invariants.\footnote{Le feuillet
50 est écrit par-dessus lui-même, et sa marge gauche au niveau de $2^{\circ}$ est
entièrement recouverte de ratures ; le détail de la construction ne se relève
pas, seule la forme de la conclusion porte.}

$3^{\circ}$ Reste la réciproque de $b)$. Il propose de la reprendre en prenant,
si $K \neq \overline{\mathbb{Q}}$, \emph{toutes} les immersions
$K \to \mathbb{C}$ à la fois, et en conjuguant au besoin avec la condition $a)$.

Ces trois lignes sont les dernières du dossier. La question $3^{\circ}$ est posée
et n'est pas reprise : le feuillet 51 s'arrête là, et il n'y a pas de suite.

\end{document}
